\documentclass[12pt,a4paper]{article}
\usepackage[utf8]{inputenc}
\usepackage[czech]{babel}
\usepackage{amsmath}
\usepackage{amsfonts}
\usepackage{amssymb}
\usepackage{geometry}
\usepackage{tcolorbox}
\usepackage{enumitem}
\usepackage{multicol}
\geometry{a4paper, margin=2cm}

\title{Příprava na čtvrtletní písemnou práci z matematiky - Teorie}
\date{}
\author{}

\begin{document}
\maketitle

Tento dokument slouží jako teoretická příprava na čtvrtletní písemnou práci. Témata pokrývají \textbf{limity posloupností}, \textbf{Kochovu vločku (nekonečné řady v geometrii)}, \textbf{aritmetické a geometrické posloupnosti (včetně parametrů)} a \textbf{slovní úlohy na finanční matematiku / exponenciální pokles}.

Příklady k procvičení najdeš v samostatném souboru \texttt{priprava\_na\_ctvrtletku\_priklady.md}.

\section{Limity posloupností}

Při výpočtu limit posloupností ($\lim_{n \to \infty} a_n$) používáme několik základních strategií podle typu zadání:

\begin{enumerate}
    \item \textbf{Racionální lomená funkce (podíl polynomů):} Vytkneme nejvyšší mocninu $n$ z čitatele i jmenovatele (nebo jednoduše podělíme každý člen nejvyšší mocninou jmenovatele).
    \begin{itemize}
        \item Pokud je nejvyšší mocnina v čitateli \textbf{menší} než ve jmenovateli, limita je $0$.
        \item Pokud jsou nejvyšší mocniny \textbf{stejné}, limita je podíl koeficientů u těchto mocnin.
        \item Pokud je nejvyšší mocnina v čitateli \textbf{větší}, limita je $\pm\infty$ (posloupnost diverguje).
    \end{itemize}
    \item \textbf{Exponenciální členy ($q^n$):} Vytkneme nejrychleji rostoucí exponenciálu (člen s největším základem). Pamatuj, že pro $|q| < 1$ platí $\lim_{n \to \infty} q^n = 0$.
    \item \textbf{Rozdíl odmocnin ($\sqrt{A} - \sqrt{B}$):} Usměrníme výraz vynásobením jedničkou ve tvaru $\frac{\sqrt{A} + \sqrt{B}}{\sqrt{A} + \sqrt{B}}$ s využitím vzorce $(x-y)(x+y) = x^2 - y^2$.
\end{enumerate}

\section{Kochova vločka (Nekonečné řady v geometrii)}

Kochova vločka je klasický fraktál. Konstrukce začíná rovnostranným trojúhelníkem o straně $a$. V každém kroku se prostřední třetina každé úsečky nahradí novým rovnostranným trojúhelníkem (směrem ven) a jeho základna se odstraní.

\begin{tcolorbox}[colback=blue!5!white,colframe=blue!75!black,title=Poznámka]
\textbf{Základní paradox:} Kochova vločka má \textbf{konečný obsah}, ale \textbf{nekonečný obvod}!
\end{tcolorbox}

\subsection{Odvození vzorců v nekonečnu ($n \to \infty$)}

Mějme výchozí rovnostranný trojúhelník o straně $a$.
\begin{itemize}
    \item \textbf{Obsah výchozího trojúhelníku ($S_0$):} $S_0 = \frac{a^2\sqrt{3}}{4}$
    \item \textbf{Obvod výchozího trojúhelníku ($O_0$):} $O_0 = 3a$
\end{itemize}

\subsubsection*{A) Obvod v nekonečnu ($O_\infty$)}
\begin{itemize}
    \item V kroku 0 máme $3$ úsečky délky $a$.
    \item V kroku 1 se každá úsečka nahradí $4$ úsečkami o délce $\frac{a}{3}$. Celkový počet úseček je $3 \cdot 4 = 12$, délka každé je $\frac{a}{3}$. Obvod $O_1 = 12 \cdot \frac{a}{3} = 4a$.
    \item V kroku $n$ máme $3 \cdot 4^n$ úseček o délce $\frac{a}{3^n}$.
    \item \textbf{Vzorec pro obvod v $n$-tém kroku:}
    \[ O_n = 3a \cdot \left(\frac{4}{3}\right)^n \]
    \item Protože kvocient $q = \frac{4}{3} > 1$, limita pro $n \to \infty$ roste do nekonečna:
    \[ \mathbf{O_\infty = \lim_{n\to\infty} O_n = \infty} \]
\end{itemize}

\subsubsection*{B) Obsah v nekonečnu ($S_\infty$)}
\begin{itemize}
    \item V 1. kroku přidáme $3$ malé trojúhelníky o straně $\frac{a}{3}$. Obsah každého z nich je $\frac{S_0}{9}$. Přidaný obsah je $3 \cdot \frac{S_0}{9} = \frac{1}{3} S_0$.
    \item Ve 2. kroku přidáme $3 \cdot 4 = 12$ trojúhelníků o straně $\frac{a}{9}$. Obsah každého je $\frac{S_0}{81}$. Přidaný obsah je $12 \cdot \frac{S_0}{81} = \frac{4}{27} S_0$.
    \item V $n$-tém kroku přidáme $3 \cdot 4^{n-1}$ trojúhelníků o straně $\frac{a}{3^n}$. Obsah každého je $\frac{S_0}{9^n}$.
    \item Celkový obsah $S_\infty$ je součtem výchozího obsahu a nekonečné geometrické řady přidaných obsahů:
    \[ S_\infty = S_0 + \sum_{n=1}^{\infty} 3 \cdot 4^{n-1} \frac{S_0}{9^n} = S_0 \left(1 + \frac{3}{9} \sum_{n=1}^{\infty} \left(\frac{4}{9}\right)^{n-1}\right) \]
    Geometrická řada uvnitř sumy má první člen $a_1 = 1$ a kvocient $q = \frac{4}{9} < 1$. Její součet je:
    \[ s = \frac{a_1}{1-q} = \frac{1}{1 - \frac{4}{9}} = \frac{9}{5} \]
    Dosadíme zpět:
    \[ S_\infty = S_0 \left(1 + \frac{1}{3} \cdot \frac{9}{5}\right) = S_0 \left(1 + \frac{3}{5}\right) = \frac{8}{5} S_0 \]
    \item Dosazením $S_0 = \frac{a^2\sqrt{3}}{4}$ získáme výsledný vzorec pro limitní obsah:
    \[ \mathbf{S_\infty = \frac{8}{5} \cdot \frac{a^2\sqrt{3}}{4} = \frac{2a^2\sqrt{3}}{5}} \]
\end{itemize}

\section{Posloupnosti a vlastnosti}

\begin{tcolorbox}[colback=red!5!white,colframe=red!75!black,title=Varování]
\textbf{Pozor na chyták u aritmetické posloupnosti!} U soustav rovnic si dávej pozor na typ posloupnosti. Pokud je zadána soustava, která pro aritmetickou posloupnost vede na nesmysl (např. $60 = 252$), pak soustava nemá řešení. Učitelé někdy dělají chybu při přepisu zadání z geometrické na aritmetickou.
\end{tcolorbox}

\subsection{Klíčová vlastnost geometrické posloupnosti}
Pro tři po sobě jdoucí členy geometrické posloupnosti $a_1, a_2, a_3$ platí klíčová vlastnost (plynoucí z rovnosti kvocientů $\frac{a_2}{a_1} = \frac{a_3}{a_2}$):
\[ a_2^2 = a_1 \cdot a_3 \]

Tohoto vztahu se hojně využívá pro určení neznámého parametru v zadání typu "Určete $x$ tak, aby daná čísla tvořila GP".

\subsection{Vlastnost aritmetické posloupnosti}
Pro tři po sobě jdoucí členy aritmetické posloupnosti zase platí:
\[ 2a_2 = a_1 + a_3 \]

\subsection{Nekonečné geometrické řady}
Geometrická řada konverguje (má konečný součet) právě tehdy, když platí podmínka $|q| < 1$. 
Pokud je tato podmínka splněna, pak součet řady je:
\[ s = \frac{a_1}{1 - q} \]

Součet prvních $n$ členů aritmetické posloupnosti je dán vzorcem:
\[ s_n = \frac{n}{2} (a_1 + a_n) \]

\section{Slovní úlohy a Finanční matematika}

U slovních úloh na úročení nebo úbytek hodnoty používáme vzorec pro složené úročení / geometrický růst (úbytek):
\[ K_n = K_0 \cdot (1 \pm i)^n \]
\begin{itemize}
    \item $K_0$ je počáteční hodnota (jistina, intenzita, atd.).
    \item $i$ je úroková sazba nebo procento změny vyjádřené jako desetinné číslo (např. $5\% = 0{,}05$).
    \item Znaménko $+$ platí pro růst (úročení), $-$ pro pokles (absorpce světla, inflace).
    \item $n$ je počet období (počet desek, let, měsíců).
\end{itemize}

\section{Přehled vzorců (Tahák)}

\begin{multicols}{2}

\subsection*{Aritmetická posloupnost}
\begin{align*}
a_n &= a_1 + (n - 1) \cdot d \\
a_r &= a_s + (r - s) \cdot d \\
s_n &= \frac{n}{2} \cdot (a_1 + a_n) \\
a_n &= \frac{a_{n-1} + a_{n+1}}{2} \\
a_{n+1} - a_n &= d
\end{align*}

\subsection*{Geometrická posloupnost}
\begin{align*}
a_n &= a_1 \cdot q^{n-1} \\
a_r &= a_s \cdot q^{r-s} \\
s_n &= a_1 \cdot \frac{q^n - 1}{q - 1} \\
a_n^2 &= a_{n-1} \cdot a_{n+1} \\
\frac{a_{n+1}}{a_n} &= q \\
a_n &= a_1 \cdot \left(1 \pm \frac{P}{100}\right)^n
\end{align*}

\subsection*{Limity}
\begin{align*}
&\lim_{n \to \infty} \left(\frac{1}{n}\right) = 0 \\
&\lim_{n \to \infty} \left(\frac{n}{1}\right) = \infty \\
&\ln x < x^k < a^x < x! < x^x \\
&\lim_{n \to \infty} \left(1 + \frac{1}{n}\right)^n = e \\
&\lim_{n \to \infty} \left(1 - \frac{1}{n}\right)^n = \frac{1}{e} \\
&\lim_{n \to \infty} \left(1 + \frac{1}{n}\right)^{n \cdot x} = e^x \\
&\lim_{n \to \infty} \left(1 + \frac{1}{x \cdot n}\right)^n = \sqrt[x]{e}
\end{align*}

\columnbreak

\subsection*{Logaritmy}
\begin{align*}
&\log_a(x) = y \implies a^y = x \\
&\log_a(1) = 0 \\
&\log_a(a) = 1 \\
&\log_a(x \cdot y) = \log_a(x) + \log_a(y) \\
&\log_a\left(\frac{x}{y}\right) = \log_a(x) - \log_a(y) \\
&\log_a(x^r) = r \cdot \log_a(x) \\
&\log_a(x) = \frac{\log_b(x)}{\log_b(a)}
\end{align*}

\subsection*{Mocniny}
\begin{align*}
&a^0 = 1 \\
&a^r \cdot a^s = a^{r+s} \\
&\frac{a^r}{a^s} = a^{r-s} \\
&(a^r)^s = a^{r \cdot s} \\
&(a \cdot b)^r = a^r \cdot b^r \\
&\left(\frac{a}{b}\right)^r = \frac{a^r}{b^r} \\
&a^{-r} = \frac{1}{a^r}
\end{align*}

\subsection*{Odmocniny}
\begin{align*}
&\sqrt[n]{a^m} = a^{\frac{m}{n}} \\
&\sqrt[n]{a \cdot b} = \sqrt[n]{a} \cdot \sqrt[n]{b} \\
&\sqrt[n]{\frac{a}{b}} = \frac{\sqrt[n]{a}}{\sqrt[n]{b}} \\
&\sqrt[m]{\sqrt[n]{a}} = \sqrt[m \cdot n]{a} \\
&\sqrt[n]{a^m} = \sqrt[n \cdot k]{a^{m \cdot k}}
\end{align*}

\end{multicols}

\section{Rychlý checklist před písemkou:}
\begin{itemize}[label=$\square$]
    \item Umím spočítat limitu podílu polynomů (porovnáním stupňů a vytknutím).
    \item Umím usměrnit limitu typu $\infty - \infty$ obsahující odmocniny.
    \item Pamatuji si, že obvod Kochovy vločky jde k $\infty$ a obsah k $\frac{8}{5} S_0$.
    \item Vím, že pro tři po sobě jdoucí členy GP platí $a_2^2 = a_1 \cdot a_3$ a umím použít substituci.
    \item Nezapomenu otočit znaménko nerovnice, pokud dělím záporným logaritmem (např. $\ln(0{,}95)$).
    \item Chápu rozdíl mezi jednoduchým úročením ($K_n = K_0(1+n \cdot i)$) a složeným úročením ($K_n = K_0(1+i)^n$).
\end{itemize}

\end{document}
